\subsection{Input Data for ``Additional Surfaces"}
\label{sec2.3b}
%\addcontentsline{toc}{section}{\protect\numberline{2.3B}{Input Data
%for "Additional Surfaces"}}
%\setcounter{section}{3}

\textbf{General remarks} 

\medskip 

Internally each additional surface is defined by an algebraic equation and some algebraic inequalities specifying the boundary of that surface, i.e., that part of the surface which is seen by the test particles.
By changing the sign of all coefficients in the algebraic equation the orientation of the surface normal vectors are reversed.

The intersection with the nearest surface along a trajectory is found by checking surfaces in the sequence of their input (in subroutine TIMEA of the geometry block GEO3D).
This must be taken into consideration, if there are two parts of one surface specified by different boundaries and with non-empty intersection.
In this case always the surface later in the input sequence is seen by the test particles.

One can define two distinct plane surfaces as one surface of second order, provided that this is compatible with the options available for surface boundaries.

We will refer to positive and negative directions of flights for particles intersecting a surface.
By ``positive" it is meant that the angle between the trajectory of the particle and the surface normal at the point of intersection is less than 90 degrees, ``negative" has the opposite meaning.

\textbf{speeding up geometry calculations:} 

For optimization of CPU performance the checking of intersections with surfaces along a test flight can be abandoned depending upon the current position of a test particle.

For particles in cell no.\ NCELL all surfaces MSURF with IGJUM3(NCELL, MSURF) = 1 are not checked.
For particles on surface MS all other surfaces MSURF with IGJUM1(MS, MSURF) = 1 are not checked.

The integer matrix IGJUM1 can be set either in this block, see CH1 cards described below, or in some user-specified problem specific routines (e.g. GEOUSR).
The integer matrix IGJUM3 can be set via the CH3 cards in input block 8, or, again, e.g.\ in problem specific routines such as GEOUSR.

\textbf{ignorable coordinates:} 

Surfaces with ignorable coordinates are specified either by setting the corresponding coefficients in the surface equation to zero, or, in case of 2-point, 3-point, 4-point of 5-point options, by setting the corresponding coordinates of the points to -1.D20 and +1.D20, respectively.
\smallskip

\textbf{coordinate systems, transformations:} 

All surfaces are specified in Cartesian coordinates.

In case of NLTRA (toroidal geometry approximation), surfaces must be defined in the local coordinate system of toroidal cells centered at the ``magnetic axis", (xa,ya) = (RMTOR,0.
), i.e.,
excluding the major radius of the torus.

To facilitate input of the geometrical data, each single (or a set of) surface(s) can be transformed by a Cartesian mapping after specification of the surface coefficients and boundaries in some ``convenient" coordinate frame, by including ``TRANSFORM cards" at the end of an input block for a particular surface.
One such ``TRANSFORM-deck" can act on an entire range of surfaces, but only on those which have been read previous to the ``TRANSFORM-deck".
Hence, in order to transform the entire set of additional surfaces by one deck, this deck must come after the last surface (no.\ NLIMI, see below).

\textbf{Particle-surface interaction models} 

Although the (local) data for particle surface interaction models for each specific additional surface can be read in this input block, their meaning is described in block 6 (\cref{sec2.6}) together with the globally valid input data for particle surface interactions.

\medskip

\textbf{The Input Block} 

\begin{lstlisting}
*** 3B. Data for Additional Surfaces
      NLIMI
C     ``CH0-cards" (may be omitted)
C     (format: 3A,69A)
\end{lstlisting} 

arbitrary number of strings $\pm n/m$, separated by blanks.
n and m
must be integer variables $1 \le n,m \le NLIMI$.

\begin{lstlisting}
      DO  ILIMI=1,NLIMI
        TXTSFL
C       ``CH1-cards" (may be omitted)
C       (format: 3A,69A)
\end{lstlisting}

arbitrary number of strings $\pm n/m$, separated by blanks. n and m
must be integer variables $1\le n,m\le NLIMI$.

\begin{lstlisting}
C       ``CH2-cards" (may be omitted)
C       (format: 3A,69A)
\end{lstlisting}

arbitrary number of strings $\pm$ n/m, separated by blanks. n and m
must be integer variables $1\le n,m\le NLIMI$.

\begin{lstlisting}
        RLBND  RLARE  RLWMN  RLWMX
        ILIIN ILSIDE ILSWCH ILEQUI ILTOR ILCOL ILFIT ILCELL 
     .    ILBOX ILPLG
        IF (RLBND.LT.0.) THEN
          A0LM  A1LM  A2LM  A3LM  A4LM  A5LM
          A6LM  A7LM  A8LM  A9LM
\end{lstlisting}

let RLBND = -KL, then first read L cards:

\begin{lstlisting}
          ALIMS  XLIMS  YLIMS  ZLIMS
\end{lstlisting}

and then read K blocks next:

\begin{lstlisting}
          ALIMS0  XLIMS1  YLIMS1  ZLIMS1  XLIMS2  YLIMS2
          ZLIMS2  XLIMS3  YLIMS3  ZLIMS3
        ENDIF
        IF (RLBND.EQ.0.) THEN
          A0LM  A1LM  A2LM  A3LM  A4LM  A5LM
          A6LM  A7LM  A8LM  A9LM
        ENDIF
        IF (RLBND.GT.0..AND.RLBND.LT.2.) THEN
          A0LM  A1LM  A2LM  A3LM  A4LM  A5LM
          A6LM  A7LM  A8LM  A9LM
          XLIMS1  YLIMS1  ZLIMS1  XLIMS2  YLIMS2  ZLIMS2
        ENDIF
        IF (RLBND.GE.2.) THEN
          P1(1,..)  P1(2,..)  P1(3,..)  P2(1,..)  P2(2,..)  
     .      P2(3,..)
          IF (K .GT. 2) THEN
            P3(1,..)  P3(2,..)  P3(3,..)  P4(1,..)  P4(2,..)  
     .        P4(3,..)
            IF (K .GT. 4) THEN
              P5(1,..)  P5(2,..)  P5(3,..)
            ENDIF
          ENDIF
        ENDIF
\end{lstlisting}

optional for non transparent surfaces (ILIIN $>$ 0):

The next 3 (or 4) lines comprise the block for local particle-surface interaction data (in analytic terminology: the boundary condition at this surface element).
If they are omitted, the default particle-surface model is activated for this particular surface element, see \cref{sec2.6}.

\begin{lstlisting}
        ILREF  ILSPT ISRS  ISRC
        ZNML EWALL EWBIN TRANSP(1,N) TRANSP(2,N) FSHEAT
        RECYCF  RECYCT  RECPRM  EXPPL  EXPEL  EXPIL
C (following line may be omitted, then: default sputter 
C  model, see \cref{sec2.6})
        RECYCS  RECYCC  SPTPRM
\end{lstlisting}

also optional for non transparent surfaces (ILIIN $>$ 0):

read a surface interaction model identifier to link one of the surface
``local reflection
models" defined in block 6 (\cref{sec2.6}) below to this current surface.
The presence of such a link is identified by the string \lstinline|SURFMOD_|, followed by a name \lstinline|modname|.
Later, in input block 6 (\cref{sec2.6}), a ``local reflection model" block with that name \lstinline|modname| must be included.
This option allows quick changes of particle-surface interaction parameters, affecting many surfaces at once, by changes in just one location of the input file.

\begin{lstlisting}
        SURFMOD_modname
\end{lstlisting}

optional: read one or several blocks of five cards each for orthogonal
mapping.
The presence of each such block is identified by the first card of that block containing the string $'TRANSFORM'$ followed by the other four cards:

\begin{lstlisting}
        TRANSFORM
        ITINI,ITEND
        XLCOR,YLCOR,ZLCOR
        XLREF,YLREF,ZLREF
        XLROT,YLROT,ZLROT,ALROT

C (end of ILIMI loop)
      ENDDO
\end{lstlisting}

\medskip

\textbf{Meaning of the Input Variables for additional surfaces} 

\medskip 

\begin{description}
  \item[CH0 n1/m1 n2/m2 \dots] 

        surfaces from the range n1 to m1, n2 to m2, \dots, are ignored by EIRENE.
        Specifying a surface in such a CH0-card is identical to taking it out from the input file.
        It may be more convenient in some cases, however, to use the CH0 option, because the labelling index of the remaining valid surfaces is not altered then.
        No input is read for these surfaces, and the input segment for the next valid surface (identified by the string '*text') is read next.
  \item[NLIMI] number of surfaces in the input block
  \item[TXTSUR] Text to identify a surface (``name of the surface") on the printout file
  \item[CH1(ILIMI) n1/m1 n2/m2 \dots] 

        surfaces from the range n1 to m1, n2 to m2, \dots, are considered invisible for a particle located on this current surface ILIMI.
        Intersection of trajectories starting from surface ILIMI with those ``invisible" surfaces is not checked.
  \item[CH2(ILIMI) n1/m1 n2/m2 \dots] 

        Only for second order surfaces ILIMI, n1-m1, \dots The first of the two possible intersections is ignored for particles located on surface no.\ ILIMI.
\end{description}

\textsl{General Data for Surfaces}
\begin{description}
  \item[RLBND] flag for different options to define the boundary of the surface
  \item[RLARE] Area (in \si{\square\centi\metre}) of the surface element which is seen by the test particles.
        (Default: 666.0) (needed only for scaling of non default surface averaged tallies) If RLARE is not specified here, (i.e., if a value less than or equal to zero is read) then EIRENE tries to evaluate this area itself.
        For some surfaces this is still not possible automatically.
  \item[RLWMN] lower weight limit for space weight window for particles crossing the surface in positive direction.
        (not in use)
  \item[RLWMX] upper weight limit for space weight window for particles crossing the surface in positive direction.
        (not in use)
  \item[ILIIN] defines the type of surface
        \begin{description}
          \item[$> 0$] non-transparent surface
                \begin{description}
                  \item[$= 1$] reflecting, partly or purely absorbing surface.

                        local reflection model has to be specified unless default model is
                        to be used; all surface tallies (see \cref{tab3})
                        are updated and a switch can be operated.
                  \item[$= 2$] purely absorbing surface ;.
                        surface tallies for incident fluxes
                        (i.e., POT\dots\ and EOT\dots\ tallies in \cref{tab3})
                        are updated and the particle history is stopped then.
                  \item[$= 3$] mirror for incident test particles.
                        I.e., specular reflection for neutral test particles, and for charged test particles the sign of the velocity component parallel to the B-field is reversed.
                  \item[$= 4$] periodicity surface, with regard to x, y, or z coordinate, depending upon whether this surface is a standard x, y, or z grid surface, respectively.
                        Move particle to x / radial grid surface no.\ m, m integer (or to y / poloidal or to z / toroidal surface no.\ m, respectively) and continue track from there with otherwise identical particle parameters.

                        This option is currently implemented only for Cartesian grids (NLSLB and NLTRZ) and for non-default standard grid surfaces only.
                \end{description}
                The periodicity options are not fully implemented yet.
                Please contact the author for the current status of your particular version.
          \item[$\le 0$] transparent surface (for example: hole in one of the other ``additional surfaces").
                Particle and energy fluxes onto and from these surfaces do not contribute to global balances.
                \begin{description}
                  \item[$= 0$] Particle history is not interrupted 

                        No surface tallies are updated, no switches can be operated.
                        Fastest option.
                  \item[$= -1$] Particle history will be stopped and restarted.
                        A switch can be operated.
                        I.e., this surface is used only for switching (see below: ILSWCH) or re-initializing the particle's track at the point of intersection.
                        No surface tallies are updated.
                  \item[$= -2$] as -1, and, additionally:

                        if a particle is crossing the surface in the positive direction,
                        (one sided-) surface tallies are updated, e.g., by
                        default: partial particle and energy currents $J^+$ (\si{\ampere})
                        and $K^+$ (\si{\watt}).
                        These are stored in the POT\dots\ and EOT\dots\ tallies of \cref{tab3}.
                        If the particle crosses the surface in the direction opposite to the surface normal, then negative partial particle and energy currents $J^-$ (\si{\ampere}) and $K^-$ (\si{\watt}) are updated.
                        These are stored in the PRF\dots\ and ERF\dots\ tallies of \cref{tab3}.

                  \item[$= -3$] Net currents (e.g.\ $J^+ - J^-$), are evaluated, and stored on the POT\dots\ and EOT\dots\ tallies (see \cref{tab3}).
                        The PRF\dots\ and ERF\dots\ tallies are empty for these surfaces.
                  \item[$\leq -4$] Not in use.
                        Currently: same as ILIIN=-2 option.
                \end{description}
        \end{description}
  \item[ILSIDE] ~
        \begin{description}
          \item[$= 0$] both sides of the surface act as described by ILIIN option (default).
          \item[$= 1$] particles incident on the surface in the negative direction will be absorbed (i.e., ILIIN = 2 option from that side).
          \item[$= 2$] particles incident on the surface in the negative direction will be killed and the message 

                ``ERROR IN ADDCOL" 

                or 

                ``ERROR IN STDCOL" 

                will be printed.
                The contribution of these particles to the particle- and energy flux balances will be called PTRASH and ETRASH respectively.
                This option should be used for geometry testing whenever the user expects particles incident only from one side.
          \item[$= 3$] particles incident on the surface in the negative direction will not see the surface, i.e., this surface acts like a (semi) transparent surface (ILIIN = 0 option) from that side.
          \item[$= -1$] as 1, but with the opposite direction of the surface normal
          \item[$= -2$] as 2, but with the opposite direction of the surface normal
          \item[$= -3$] as 3, but with the opposite direction of the surface normal
        \end{description}
  \item[ILSWCH] = IJKLMN, i.e.\ six digits I, J, K, L, M and N
        \begin{description}
          \item[$= 0$] no switch is operated
          \item[N] EIRENE flag ITIME
                \begin{description}
                  \item[$= 1$] The calculation of the step sizes in the standard mesh is abandoned for a particle which crosses the surface in the positive direction, and is reactivated, if the particle strikes in the negative direction
                  \item[$= 2$] as 1, but with the direction of the surface normal reversed for this option.
                \end{description}
          \item[M] EIRENE flag IFPATH
                \begin{description}
                  \item[$= 1$] Abandon the calculation of the collision rates (entry into the vacuum) for a particle which is striking the surface in the direction of the surface normal.
                        For particles incident from the other direction, evaluation of collision rates is reactivated.
                  \item[$= 2$] as 1, but with the direction of the surface normal reversed.
                \end{description}
          \item[L] EIRENE flag IUPDTE
                \begin{description}
                  \item[$= 1$] Abandon the updating of volume-averaged tallies for a particle which is striking the surface in the direction of the surface normal.
                        For particles incident from the other direction, updating of volume averaged tallies is reactivated.
                  \item[$= 2$] as 1, but with the direction of the surface normal reversed for this option.
                \end{description}
        \end{description}

        I,J,K flags for switching cell numbers at transition into a different mesh cell.
        \begin{description}
          \item[K] for particles in an additional cell, i.e., not in one of the ``standard mesh" blocks:
                \begin{description}
                  \item[$= 1$] Increase the actual additional cell number NACELL for a particle striking the surface in the direction of the surface normal by ILACLL.
                        Decrease NACELL by ILACLL if the particle is striking in the negative direction.
                        Specification of ILACLL is via the input variable ILCELL, see below.
                  \item[$= 2$] as K $= 1$, but with the direction of the surface normal reversed for this option.
                \end{description}
                for particles inside the ``standard mesh ", i.e., not in the ``additional
                cell region"
                \begin{description}
                  \item[$= 1$] Increase the standard mesh block number NBLOCK for a particle striking the surface in the direction of the surface normal by ILBLCK.
                        Decrease NBLOCK by ILBLCK if the particle is striking in the negative direction.
                        Specification of ILBLCK is via the input variable ILCELL, see below.
                  \item[$= 2$] as K $= 1$, but with the direction of the surface normal reversed.
                \end{description}
          \item[J] for particles at the boundary between ``additional" and ``standard" mesh regions.
                \begin{description}
                  \item[$= 1$] entrance into standard mesh, block no.

                        NBLOCK = ILBLCK 

                        or exit from standard mesh into additional cell 

                        NACELL = ILACLL.

                        Specification of ILACLL and ILBLCK is via the input variable ILCELL, see below.
                        If ILACLL = 0, then no switch to additional cell is operated.
                        (E.g.: for surfaces which are reflecting from this side).
                  \item[$= 2$] as J $= 1$.
                        The direction of the surface normal does not matter here.
                \end{description}
          \item[I] similar to J-flag, i.e., for transitions between standard and additional meshes, but different cell number switching.
                \begin{description}
                  \item[$= 1$] Entrance into standard mesh, block no.

                        NBLOCK = NACELL+ILBLCK, 

                        if the particle is striking in the positive direction, or 

                        NBLOCK = NACELL-ILBLCK, 

                        if the particle is striking in the negative direction.
                        Exit from standard mesh into additional cell 

                        NACELL = NBLOCK+ILACLL, 

                        for a particle striking the surface in the positive direction, or 

                        NACELL = NBLOCK-ILACLL, 

                        if the particle is striking in the negative direction.

                        Specification of ILACLL and ILBLCK is via the input variable ILCELL, see below.
                  \item[$= 2$] as I = 1, but with the direction of the surface normal reversed.
                \end{description}
        \end{description}

        If a test particle history starts from a surface (NLSRF option), then ILSWCH acts as if this particle had struck the surface prior to the birth process in the positive direction.
        This default setting is only available for ILSIDE $\neq$ 0 and can (must) be overruled by the SORIFL flag (see \cref{sec2.7}), e.g.\ if a surface source needs to be defined on a surface with ILSIDE = 0.
  \item[ILEQUI] The algebraic equations for the surfaces J and IABS(ILEQUI(J)) will be described by exactly the same coefficients (up to a common sign, if ILEQUI(J) .lt.\ 0).
        For example a triangle can be specified by the three corners and another part of the same plane surface can be specified directly by its algebraic coefficients.
        To avoid round-off errors one should use the ILEQUI option in such cases, in particular if surface J is a transparent ``hole" in surface ILEQUI(J), or vice versa.

        Default: ILEQUI = 0
  \item[ILTOR] For NLTRA option only (see block 2c):

        if ILTOR$>$0 :

        the surface is defined with respect to the local coordinate system
        of the toroidal segment with ``cell-number" ILTOR, hence: $1 \le
        $ILTOR $\le $NTTRAM.

        if ILTOR=0 :

        the surface is defined with respect to any local coordinate system.
        I.e., the surface equations are taken to be the same in each local system.

        If the surface equations are z-independent, then this surface is toroidally symmetric (within the NLTRA-approximation).

        Otherwise the surface has NTTRAM-fold periodicity.

        This flag is irrelevant for NLTRZ (i.e., if cylindrical coordinates are used) or for NLTRT.

        Default: ILTOR = 0
  \item[ILCOL] Flag for the color that is used for plotting this surface on 2d or 3d geometry plots.
        If ILCOL $\le 0$  than -ILCOL is used and the surface area is filled in by that color on the 3d geometry plots.

        Default: ILCOL=1
  \item[ILFIT] = MN 

        This option is relevant only for surfaces with one ignorable coordinate, i.e.\ it only works for the $2.
          \le RLBND \le 3$. surface boundary options.

        It is a tool to facilitate a neat fitting of surfaces, in particular for connecting curved and plane surfaces (i.e.\ to avoid particle leakage due to numerical round-off errors in the algebraic surface coefficients.

        M and N (3 digits each, M may be omitted if not needed) are the numbers of the surfaces (which must have the same ignorable coordinate) the boundaries of which should match to those of the actual surface J.
        The boundaries of these neighboring surfaces M and N must be specified by the RLBND = 1 or RLBND = 1.5 option.

        The fitting is achieved by a small automatic internal modification of the surface data P1 and/or P2 of surface number J in subroutine SETFIT.
        Printout of the modifications made there is activated with the TRCSUR flag (input block 11).
  \item[ILCELL] Parameter ILBLCK and ILACLL for the ILSWCH flags described above.
        Let ILCELL = NM, with N and M being integers with 3 digits each.
        Then N = ILBLCK and M = ILACLL.

  \item[ILBOX] to be written
  \item[ILPLG] EIRENE can write out information for a finite element mesh generator to produce a grid of triangles for a multiply connected 2D domain with cracks and holes.
        The various (inner and outer) boundaries are given as polygonal lines, which are composed of selected standard grid surface segments (NLPLG option) and/or additional surfaces (2 $\le$ RLBND $<$ 3 option).

        This flag identifies closed polygonal lines composed of additional surfaces given by the 2-point option and/or of standard surfaces in the x-y-plane.
        For example if ILPLG(I) = NN, for surfaces I = I1, I2, \dots, IN, (NN a positive integer) then these IN surfaces form a closed polygonal line in the x-y-plane.
        The region inside this closed line is part of the computational domain.
        By a negative integer value of NN a closed polygonal region can be excluded from the computational domain, i.e., a hole in the domain is specified by these surfaces.
        EIRENE writes an output file appropriate for a finite element mesh generator (available from FZ-J\"ulich) to produce a triangular discretization of the resulting (possibly multiply connected) domain.
        This option can be used to discretize arbitrarily complex 2D domains with internal and external boundaries given by the additional or non-default standard surfaces.

        These finite element grids can be combined with the regular grids by using the problem specific geometry routines (see section 3) or the code interfacing routines INFCOP (see section 4).
\end{description}

\medskip

\textbf{Surface coefficients and boundaries of surfaces} 

\medskip 

The surface equation:
\begin{align}
  A0LM & + A1LM  \cdot  x + A2LM \cdot y + A3LM  \cdot z +  \nonumber      \\
       & + A4LM  \cdot x^2 + A5LM  \cdot y^2 + A6LM  \cdot z^2 + \nonumber \\
       & + A7LM  \cdot xy + A8LM  \cdot xz + A9LM  \cdot yz = 0
\end{align}
The positive surface normal $(n_x, n_y, n_z)$ depends upon the point of impact (x,y,z) and is defined by the vector
\begin{align}
  n_x & = A1LM  + 2  \cdot   A4LM  \cdot  x    + A7LM  \cdot  y    +      A8LM  \cdot  z    \nonumber \\
  n_y & = A2LM    +
  A7LM  \cdot  x    +     2  \cdot   A5LM  \cdot  y    +
  A9LM  \cdot  z   \nonumber                                                                          \\
  n_z & = A3LM    +
  A8LM  \cdot  x +  A9LM  \cdot  y +
  2  \cdot   A6LM  \cdot  z       \nonumber
\end{align}
For a plane surface the following reduction is valid:

($n_x, n_y, n_z$) = (A1LM,A2LM,A3LM) 

\medskip 

The boundary of the valid part of the surface may be described by 4 different options, depending upon the value of the flag RLBND.
\begin{description}
  \item[RLBND = 0] ~ 

        No boundary inequalities specified, i.e.\ the whole surface is seen by the test particles.
  \item[0  $<$   RLBND  $<$   2] ~ 

        \begin{description}
          \item[$= 1$] Only that part of the surface, which lies inside the right parallelepiped defined by the two vectors (XLIMS1, YLIMS1, ZLIMS1) and (XLIMS2, YLIMS2, ZLIMS2), is seen by the particles.
                I.e.\ the three inequalities 

                $XLIMS1 \le x \le XLIMS2$ 

                $YLIMS1 \le y \le YLIMS2$ 

                $ZLIMS1 \le z \le ZLIMS2$ 

                are checked at the point of intersection (x,y,z).
          \item[$= 1.5$] Complement to RLBND = 1.

                Only the surface element outside the parallelepiped is seen by the particles.
        \end{description}
  \item[RLBND $\ge$ 2] ~ 

        In this case the surface will be defined by the input of the coordinates of at least 2 and at highest 5 points on a plane surface.
        If there are only 2 points, the surface is parallel to one axis.
        If there are 3 or more points, then the boundary of this plane surface is a closed polygon $(P_1, \dots, P_n, P_1)$.
        Therefore, the correct order of points at input is relevant.
        The orientation of the positive surface normal vector is defined by the first 3 points, and it is given by the vector product $(P3 - P1) \times (P3 - P2)$.
        Thus, the orientation can be reversed e.g.\ by interchanging $P_2$ and $P_3$.
        \begin{description}
          \item[$= 2.1$] plane surface parallel to z axis.
                The surface equation of this plane reads ax + by + c = 0 with the coefficients a,b and c such that the points $P_1, P_2$ lie on this surface and the valid part of that surface ranges from $P_1$ to $P_2$ in the xy-plane.
                The z-coordinates of these two points define the boundaries in z direction
          \item[$= 2.2$] Complement to RLBND = 2.1
          \item[$= 2.4$] as RLBND=2.1 option, but with z and y exchanged.
                I.e., now the y coordinates of the points $P_1, P_2$ are the boundaries of the surface ax+bz+c=0 in y direction.
          \item[$= 2.5$] Complement to RLBND = 2.4
          \item[$= 2.7$] as RLBND=2.1 option, but with z and x exchanged.
                I.e., now the x coordinates of the points $P_1, P_2$ are the boundaries of the surface ay+bz+c+0 in x direction.
          \item[$= 2.8$] Complement to RLBND = 2.7
          \item[$= 3$] plane triangle defined by the corners $P_1, P_2, P_3$ 

                $P_1$=(P1(1),P1(2),P1(3)) 

                $P_2$=(P2(1),P2(2),P2(3)) 

                $P_3$=(P3(1),P3(2),P3(3))
          \item[$= 3.5$] complement to RLBND = 3; only The plane surface outside the triangle is seen by the test particles.
          \item[$= 4$] plane quadrangle; surface inside the polygon 

                ($P_1, P_2, P_4, P_3, P_1$).

                Here $P_1, P_2, P_3$ are as in the RLBND=3 option, and $P_4$ = (P4(1), P4(2), P4(3)) 

                Thus this surface is the union of the triangles with vertices $P_1$, $P_2$, $P_3$ and $P_2$, $P_4$, $P_3$ respectively.
          \item[$= 4.5$] complement to RLBND = 4; only the part of the plane surface outside the quadrangle is seen by the test particles
          \item[$= 5$] plane quint-angle; surface inside the polygon ($P_1$, $P_2$, $P_4$, $P_5$, $P_3$, $P_1$) $P_1$, $P_2$, $P_3$, $P_4$ as RLBND=4, and $P_5$ = (P5(1), P5(2), P5(3))
          \item[$= 5.5$] complement to RLBND = 5; only the part of the plane surface outside the quint-angle is seen by the test particles.
        \end{description}
  \item[RLBND < 0] ~ 

        -KL 

        The surface is bounded by L linear inequalities and by K second order inequalities.

        $ALIMS + XLIMS  *  x + YLIMS  *  y + ZLIMS  *  z \le 0 $ \raggedleft (L inequalities) 

        \begin{align}
          ALIMS0 & + XLIMS1  \cdot  x + YLIMS1  \cdot  y + ZLIMS1  \cdot  z \nonumber \\
                 & + XLIMS2  \cdot  x^2 + YLIMS2  \cdot  y^2 + ZLIMS2  \cdot  z^2
          \nonumber                                                                   \\
                 & + XLIMS3  \cdot  xy + YLIMS3  \cdot  xz + ZLIMS3  \cdot  yz \le 0
          \nonumber
        \end{align}
        \raggedleft (K inequalities)
\end{description}

The meanings of the input variables for the local reflection model are described below (see: \cref{sec2.6}: ``Input Data for Surface Interaction Models").

The meanings of the input variables for a transformation block are:

\begin{description}
  \item[ITINI, ITEND] ~ 

        The transformation defined by the next 3 cards is carried out for additional surfaces from number ITINI through ITEND The transformation is carried out as soon as this transformation deck is found.
        Hence, if such a transformation deck is found after ``additional surface" no.\ IS, then one must guarantee: ITINI $\le$ IS and ITEND $\le$ IS.

  \item[XLCOR, YLCOR, ZLCOR] ~ 

        Translation:

        The origin of the coordinate system is shifted by the vector 

        XLCOR,YLCOR,ZLCOR
  \item[XLREF, YLREF, ZLREF] ~ 

        Reflection:

        to be written

        No transformation if XLREF,YLREF,ZLREF = (0,0,0)
  \item[XLREF, YLREF, ZLREF] ~ 

        Rotation:

        The vector XLROT,YLROT,ZLROT defines the axis of rotation.

        ALROT is the angle of rotation (in degrees)

        No transformation if ALROT=0 or axis of rotation= (0,0,0)
\end{description}
